% GENERATED FILE. Edit canonical lessons, then run npm run generate:books.
% canonical-source-hash: fnv1a64:f2bcfef0a7a0d0a8
% canonical-lessons: BN-C78-rat, BN-C78-sondhya, BN-C78-sokal, BN-C78-soptaho, BN-C78-mash

\chapter{Night, Evening, Morning, Week, Month}
\label{ch:bn-night-evening-morning-week-month}
\hlchaptermodality{78}

\begin{chapteropening}
\textbf{By the end of this chapter:} \emph{I can say night, evening, morning, a week and a month.}

Complete the last lesson of chapter 78: I can say night, evening, morning, a week and a month.
\end{chapteropening}

\section[\texorpdfstring{\bn{রাত} (rāt)}{রাত (rāt)}]{\bn{রাত} (rāt) — night}

\label{lesson:BN-C78-rat}



\begin{quote}
Before the new one: say the Bengali for beautiful, then the Bengali for a day.
\end{quote}

\subsection*{You'll want to know: \bn{রাত}}
\textbf{\bn{রাত}} — \emph{rāt} — \textquotedblleft{}night\textquotedblright{}.

\begin{cousinweb}
\textbf{\bn{রাত}} (\emph{rāt}) is a tadbhava, a word inherited from Sanskrit and reshaped by sound change. It comes from \emph{rātrī}, \textquotedblleft{}night\textquotedblright{} (CDIAL 10702).
\end{cousinweb}

\begin{grammarlens}[title={What you've built}]
One more word for this chapter.
\end{grammarlens}

\subsection*{Your turn}
\begin{itemize}
\raggedright
  \item \emph{Say it:} \emph{rāt}
  \item \emph{Say it:} \emph{rāt}, once more
  \item \emph{Say it:} say \emph{rāt}
  \item \emph{Recall:} say the Bengali for beautiful, then the Bengali for a day, then say \emph{rāt} again
\end{itemize}

\begin{tcolorbox}[breakable,colback=teal!4,colframe=teal!35!black,title={Before you move on}]
What does \bn{রাত} mean? (\textquotedblleft{}Night\textquotedblright{}.) Say it once more.
\end{tcolorbox}
\section[\texorpdfstring{\bn{সন্ধ্যা} (sôndhyā)}{সন্ধ্যা (sôndhyā)}]{\bn{সন্ধ্যা} (sôndhyā) — evening}

\label{lesson:BN-C78-sondhya}



\begin{quote}
Before the new one: say the Bengali for a day, then the Bengali for night.
\end{quote}

\subsection*{You'll want to know: \bn{সন্ধ্যা}}
\textbf{\bn{সন্ধ্যা}} — \emph{sôndhyā} — \textquotedblleft{}evening\textquotedblright{}.

\begin{cousinweb}
\textbf{\bn{সন্ধ্যা}} (\emph{sôndhyā}) is Sanskrit \emph{saṁdhyā}, \textquotedblleft{}twilight\textquotedblright{}, taken over whole as a learned loan, a tatsama (CDIAL 12918). The word inherited from it is \emph{sāṁjh}, \textquotedblleft{}evening\textquotedblright{}, which CDIAL lists.
\end{cousinweb}

\begin{grammarlens}[title={What you've built}]
One more word for this chapter.
\end{grammarlens}

\subsection*{Your turn}
\begin{itemize}
\raggedright
  \item \emph{Say it:} \emph{sôndhyā}
  \item \emph{Say it:} \emph{sôndhyā}, once more
  \item \emph{Say it:} say \emph{sôndhyā}
  \item \emph{Recall:} say the Bengali for a day, then the Bengali for night, then say \emph{sôndhyā} again
\end{itemize}

\begin{tcolorbox}[breakable,colback=teal!4,colframe=teal!35!black,title={Before you move on}]
What does \bn{সন্ধ্যা} mean? (\textquotedblleft{}Evening\textquotedblright{}.) Say it once more.
\end{tcolorbox}
\section[\texorpdfstring{\bn{সকাল} (sôkāl)}{সকাল (sôkāl)}]{\bn{সকাল} (sôkāl) — morning}

\label{lesson:BN-C78-sokal}



\begin{quote}
Before the new one: say the Bengali for night, then the Bengali for evening.
\end{quote}

\subsection*{You'll want to know: \bn{সকাল}}
\textbf{\bn{সকাল}} — \emph{sôkāl} — \textquotedblleft{}morning\textquotedblright{}.

\begin{cousinweb}
\textbf{\bn{সকাল}} (\emph{sôkāl}) matches Sanskrit \emph{sakāla}, \textquotedblleft{}seasonable\textquotedblright{}, whose adverb \emph{sakālam} means \textquotedblleft{}early in the morning\textquotedblright{}. CDIAL has that entry but lists no Bengali form under it (CDIAL 13067).
\end{cousinweb}

\begin{grammarlens}[title={What you've built}]
One more word for this chapter.
\end{grammarlens}

\subsection*{Your turn}
\begin{itemize}
\raggedright
  \item \emph{Say it:} \emph{sôkāl}
  \item \emph{Say it:} \emph{sôkāl}, once more
  \item \emph{Say it:} say \emph{sôkāl}
  \item \emph{Recall:} say the Bengali for night, then the Bengali for evening, then say \emph{sôkāl} again
\end{itemize}

\begin{tcolorbox}[breakable,colback=teal!4,colframe=teal!35!black,title={Before you move on}]
What does \bn{সকাল} mean? (\textquotedblleft{}Morning\textquotedblright{}.) Say it once more.
\end{tcolorbox}
\section[\texorpdfstring{\bn{সপ্তাহ} (sôptāhô)}{সপ্তাহ (sôptāhô)}]{\bn{সপ্তাহ} (sôptāhô) — a week}

\label{lesson:BN-C78-soptaho}



\begin{quote}
Before the new one: say the Bengali for evening, then the Bengali for morning.
\end{quote}

\subsection*{You'll want to know: \bn{সপ্তাহ}}
\textbf{\bn{সপ্তাহ}} — \emph{sôptāhô} — \textquotedblleft{}a week\textquotedblright{}.

\begin{cousinweb}
\textbf{\bn{সপ্তাহ}} (\emph{sôptāhô}) is Sanskrit \emph{saptāha}, \textquotedblleft{}period of seven days\textquotedblright{}, taken over whole as a learned loan, a tatsama. It is \emph{sapta}, \textquotedblleft{}seven\textquotedblright{}, + \emph{aha}, \textquotedblleft{}day\textquotedblright{}.
\end{cousinweb}

\begin{grammarlens}[title={What you've built}]
One more word for this chapter.
\end{grammarlens}

\subsection*{Your turn}
\begin{itemize}
\raggedright
  \item \emph{Say it:} \emph{sôptāhô}
  \item \emph{Say it:} \emph{sôptāhô}, once more
  \item \emph{Say it:} say \emph{sôptāhô}
  \item \emph{Recall:} say the Bengali for evening, then the Bengali for morning, then say \emph{sôptāhô} again
\end{itemize}

\begin{tcolorbox}[breakable,colback=teal!4,colframe=teal!35!black,title={Before you move on}]
What does \bn{সপ্তাহ} mean? (\textquotedblleft{}A week\textquotedblright{}.) Say it once more.
\end{tcolorbox}
\section[\texorpdfstring{\bn{মাস} (māsh)}{মাস (māsh)}]{\bn{মাস} (māsh) — a month}

\label{lesson:BN-C78-mash}



\begin{quote}
Before the new one: say the Bengali for morning, then the Bengali for a week.
\end{quote}

\subsection*{You'll want to know: \bn{মাস}}
\textbf{\bn{মাস}} — \emph{māsh} — \textquotedblleft{}a month\textquotedblright{}.

\begin{cousinweb}
\textbf{\bn{মাস}} (\emph{māsh}) is inherited from Sanskrit \emph{māsa}, \textquotedblleft{}moon, month\textquotedblright{}, and has come down with almost no change (CDIAL 10104).
\end{cousinweb}

\begin{grammarlens}[title={What you've built}]
One more word for this chapter.
\end{grammarlens}

\subsection*{Your turn}
\begin{itemize}
\raggedright
  \item \emph{Say it:} \emph{māsh}
  \item \emph{Say it:} \emph{māsh}, once more
  \item \emph{Say it:} say \emph{māsh}
  \item \emph{Recall:} say the Bengali for morning, then the Bengali for a week, then say \emph{māsh} again
\end{itemize}

\begin{tcolorbox}[breakable,colback=teal!4,colframe=teal!35!black,title={Before you move on}]
What does \bn{মাস} mean? (\textquotedblleft{}A month\textquotedblright{}.) Say it once more.
\end{tcolorbox}
